\rtmtight
\begin{tblr}{width=\textwidth,colspec={|Q[c,m,2.1cm]|X[l,m]|},hlines,vlines,hspan=minimal,rowsep=1.5pt,colsep=3pt,row{1}={bg=headgray,font=\bfseries}}
\SetCell{c,m}\hfil\logo\hfil & \textbf{POLITEKNIK NEGERI MALANG}\par\textbf{JURUSAN TEKNOLOGI INFORMASI}\par\textbf{PROGRAM STUDI: D4 TEKNIK INFORMATIKA} \\
\SetCell[c=2]{l} \textbf{RENCANA TUGAS MAHASISWA (RTM) DAN RUBRIK PENILAIAN} & \\
Nama Mata Kuliah & Kewarganegaraan \\
Kode Mata Kuliah & RTI253003 \quad \textbf{sks / Jam perminggu:} 2 SKS/3 jam \quad \textbf{Semester:} 3 \\
Dosen Pengampu & \begin{enumerate}\item Dr. Widaningsih Condrowardani, S.Psi, SH, MH\end{enumerate} \\
\end{tblr}
\assessmentblock{Tugas Analisis Kewarganegaraan}{Tugas terstruktur (analisis kasus dan kajian teori)}{SCPMK0102-02101, SCPMK0102-02102, SCPMK0102-02103, SCPMK0102-02104}{Mahasiswa menyelesaikan tugas analisis kewarganegaraan setiap pertemuan (Mgg 1--6, 9--13) yang mencakup identitas nasional, konstitusi, hubungan negara-warga negara, negara hukum, demokrasi, HAM, dan wawasan nusantara sesuai materi yang sedang dipelajari.}{Mengkaji teori dan konsep, menganalisis studi kasus atau fenomena aktual terkait materi, mengaitkan dengan nilai-nilai Pancasila, dan mendokumentasikan hasil analisis dalam bentuk tulisan terstruktur.}{Laporan analisis tertulis, makalah, atau lembar kerja sesuai format yang ditentukan dosen.}{Ketepatan pemahaman dan penerapan konsep kewarganegaraan & 9 & Konsep kewarganegaraan (identitas nasional, demokrasi, HAM, wawasan nusantara) dipahami secara tepat dan diaplikasikan pada analisis kasus dengan argumentasi yang kuat. & Pemahaman konsep sebagian besar benar, tetapi penerapan pada kasus masih kurang mendalam. & Konsep tidak dipahami dengan benar atau tidak dapat diterapkan pada kasus yang dianalisis. \\ \\
Kualitas analisis dan argumentasi & 8 & Analisis disertai argumen logis, bukti atau contoh relevan, dan dikaitkan dengan konteks kebangsaan dan nilai Pancasila. & Analisis cukup logis, tetapi argumen atau keterkaitannya dengan konteks kebangsaan masih kurang kuat. & Analisis tidak logis, tidak berargumen, atau tidak berkaitan dengan konteks kewarganegaraan. \\ \\
Kerapian dan sistematika tulisan & 3 & Tulisan sistematis, bahasa komunikatif, dan dokumentasi lengkap sesuai format. & Tulisan cukup sistematis tetapi kurang rapi atau ada format yang belum terpenuhi. & Tulisan tidak sistematis atau tidak sesuai format yang ditentukan. \\}

\assessmentblock{Kuis Konsep Kewarganegaraan}{Kuis tertulis}{SCPMK0102-02101, SCPMK0102-02102}{Mahasiswa mengerjakan kuis konsep kewarganegaraan pada minggu ke-7 yang mencakup materi identitas nasional, negara dan konstitusi, hubungan negara dan warga negara, serta negara hukum (materi Mgg 1--6) secara mandiri dan jujur.}{Mengerjakan soal pilihan ganda dan/atau esai pendek secara individu dengan sifat closed book sesuai jadwal asesmen.}{Lembar jawaban kuis.}{Penguasaan konsep identitas nasional dan kenegaraan & 6 & Jawaban tepat dan lengkap; pengertian identitas nasional, konstitusi, hak dan kewajiban warga negara, serta ciri negara hukum dijelaskan benar sesuai soal. & Sebagian besar jawaban benar, tetapi ada kekeliruan dalam penjelasan konsep tertentu. & Banyak jawaban tidak sesuai atau konsep kewarganegaraan tidak dikuasai. \\ \\
Kemampuan menganalisis kasus kewarganegaraan & 4 & Analisis kasus pada soal tepat, menunjukkan pemahaman tentang prinsip-prinsip kewarganegaraan dan relevansinya. & Analisis kasus sebagian benar, tetapi penerapan prinsip masih kurang tepat. & Tidak mampu menganalisis kasus atau jawaban tidak relevan dengan materi. \\}

\assessmentblock{UTS Kewarganegaraan}{Ujian tengah semester (esai dan/atau pilihan ganda)}{SCPMK0102-02101, SCPMK0102-02102}{Mahasiswa mengerjakan soal UTS yang mencakup materi minggu 1--7: identitas nasional, negara dan konstitusi, hubungan negara dan warga negara, negara hukum, dan demokrasi secara mandiri dan jujur.}{Mengerjakan soal esai dan/atau pilihan ganda secara individu dengan sifat closed book sesuai jadwal asesmen.}{Lembar jawaban UTS.}{Penguasaan konsep identitas nasional dan kenegaraan & 8 & Pengertian identitas nasional, konstitusi, hak-kewajiban warga negara, dan negara hukum dijelaskan tepat dan komprehensif sesuai soal. & Sebagian besar konsep dipahami, tetapi penjelasan pada soal tertentu kurang lengkap. & Banyak konsep tidak dikuasai atau jawaban tidak relevan dengan soal. \\ \\
Penguasaan konsep demokrasi dan negara hukum & 8 & Prinsip-prinsip demokrasi, perjalanan demokrasi Indonesia, dan ciri negara hukum dijelaskan benar dan dikaitkan dengan konteks kehidupan berbangsa. & Konsep demokrasi sebagian besar benar, tetapi keterkaitan dengan konteks kehidupan berbangsa kurang kuat. & Konsep demokrasi atau negara hukum tidak dipahami dengan benar. \\ \\
Kejelasan dan sistematika jawaban esai & 4 & Jawaban esai terstruktur, berargumen dengan logis, dan menggunakan bahasa yang komunikatif. & Jawaban cukup jelas tetapi kurang sistematis atau argumennya tidak kuat. & Jawaban tidak terstruktur atau tidak dapat dipahami. \\}

\assessmentblock{Proyek Berbasis Masalah Kewarganegaraan}{Project Based Learning (PBL)}{SCPMK0102-02102, SCPMK0102-02103, SCPMK0102-02104}{Mahasiswa mengerjakan proyek berbasis masalah secara berkelompok (Mgg 14--16) yang mencakup analisis implementasi wawasan nusantara, ketahanan nasional, dan integrasi nasional dalam konteks permasalahan aktual di Indonesia, kemudian mempresentasikan hasilnya.}{Mengidentifikasi permasalahan aktual terkait wawasan nusantara/ketahanan nasional/integrasi nasional, menganalisis dengan konsep kewarganegaraan, merumuskan rekomendasi, menyusun laporan proyek, dan mempresentasikan hasil kepada kelas.}{Laporan proyek kelompok, bahan presentasi (slide), dan dokumentasi diskusi kelompok.}{Ketepatan identifikasi masalah dan penerapan konsep kewarganegaraan & 6 & Masalah aktual diidentifikasi dengan tepat dan dianalisis menggunakan konsep wawasan nusantara, ketahanan nasional, atau integrasi nasional secara mendalam. & Masalah teridentifikasi, tetapi analisis menggunakan konsep kewarganegaraan masih kurang mendalam. & Masalah tidak relevan atau tidak dianalisis menggunakan konsep kewarganegaraan. \\ \\
Kualitas rekomendasi dan solusi & 5 & Rekomendasi spesifik, relevan dengan konteks kebangsaan, dapat ditindaklanjuti, dan dikaitkan dengan nilai-nilai Pancasila. & Rekomendasi cukup relevan tetapi masih umum atau tidak dikaitkan dengan nilai-nilai Pancasila. & Rekomendasi tidak relevan, tidak operasional, atau tidak menjawab masalah yang diidentifikasi. \\ \\
Kualitas presentasi dan kerjasama tim & 4 & Presentasi jelas, terstruktur, seluruh anggota tim terlibat aktif, dan mampu menjawab pertanyaan dengan baik. & Presentasi cukup jelas tetapi kurang terstruktur atau tidak semua anggota terlibat aktif. & Presentasi tidak jelas, tidak terstruktur, atau tidak mampu menjawab pertanyaan. \\}

\assessmentblock{UAS Kewarganegaraan}{Ujian akhir semester komprehensif (esai dan/atau pilihan ganda)}{SCPMK0102-02101, SCPMK0102-02102, SCPMK0102-02103, SCPMK0102-02104}{Mahasiswa mengerjakan soal UAS yang mencakup seluruh materi minggu 1--16 secara komprehensif: identitas nasional, negara dan konstitusi, hubungan negara-warga negara, negara hukum, demokrasi, HAM, wawasan nusantara, ketahanan nasional, dan integrasi nasional.}{Mengerjakan soal esai dan/atau pilihan ganda secara individu dengan sifat closed book sesuai jadwal asesmen.}{Lembar jawaban UAS.}{Penguasaan konsep kewarganegaraan secara komprehensif & 15 & Seluruh konsep kewarganegaraan (identitas nasional, demokrasi, HAM, wawasan nusantara, ketahanan, integrasi nasional) dikuasai dan dijawab tepat sesuai soal. & Sebagian besar konsep dikuasai, tetapi ada beberapa materi yang kurang dipahami secara mendalam. & Banyak konsep tidak dikuasai atau jawaban tidak sesuai dengan materi kewarganegaraan. \\ \\
Kemampuan analisis dan sintesis antar topik & 12 & Mampu mengaitkan dan mensintesis konsep-konsep kewarganegaraan dari berbagai topik untuk menjawab soal secara terpadu. & Mampu menganalisis topik tertentu, tetapi sintesis antartopik masih terbatas. & Tidak mampu menganalisis atau mensintesis konsep antar topik kewarganegaraan. \\ \\
Kejelasan, sistematika, dan bahasa jawaban & 8 & Jawaban terstruktur dengan baik, argumentatif, menggunakan bahasa Indonesia yang baik, dan mudah dipahami. & Jawaban cukup jelas tetapi sistematika atau penggunaan bahasa masih kurang baik. & Jawaban tidak terstruktur, tidak argumentatif, atau sulit dipahami. \\}

\begin{tblr}{width=\textwidth,colspec={|Q[l,m,3.2cm]|X[l,m]|},hlines,vlines,hspan=minimal,row{1-3}={bg=headgray},rowsep=1.5pt,colsep=3pt}
Jadwal Pelaksanaan: & Minggu 1--17 sesuai RPS. Setiap evaluasi memiliki RTM dan rubrik multi-kriteria. \\
Lain-Lain Yang Diperlukan: & LMS, bahan referensi kewarganegaraan, perangkat presentasi. \\
Pustaka: & Mengacu pustaka utama dan pendukung pada bagian RPS. \\
\end{tblr}
